% ch04b.tex —— Xi Yin, Aspects of Two-Dimensional Conformal Field Theories（PoS TASI2017）中译
% 覆盖 PDF 第 27–33 页（原书 p.26–32）：4.5 节续文、4.6 超共形对称性、4.7 Calabi–Yau 模型、
% 4.8 精确边缘形变、4.9 RG 流。另补译 PDF 第 34 页顶部至第 4 节结束的文字（含 (4.46)），使第 4 节完整收束。
% 本区间无图。公式编号沿用原文 (4.28)–(4.46)，文献 [N] 一律对应参考文献键 rN。译者：tr04b。

例如，领头阶条件 $R_{ij}=0$ 等价于 NLSM 单圈 beta 函数的消失，它要求靶空间 Ricci 平坦。这给靶空间加上了很不平凡的拓扑限制。一类允许 Ricci 平坦度规的偶数维紧致流形是 Calabi--Yau 流形。可以问：给定 Ricci 平坦度规为 $G_{ij}(X)$ 的 Calabi--Yau 流形 $M$，能否开启膨胀子剖面 $\Phi(X)$ 并形变度规，使 \eqref{eq:4-27} 在全阶都被解出？我不知道这个问题的答案；不过，这个问题的超共形版本答案是肯定的，下一小节就来讨论。

% [Tutor Note: 这里提到的 (4.27) 是上一小节给出的含膨胀子的 NLSM $\beta$ 函数方程（已含二圈项）；其领头阶只留下 $R_{ij}+2\nabla_i\nabla_j\Phi$，故只保留度规项时给出 $R_{ij}=0$。作者强调“全阶”可解，是因为 $\alpha'$ 修正会引入高阶曲率项。]

\subsection{超共形对称性}

带超对称的 CFT 妙处很多：其中之一是这类理论大量存在，而由于各种超对称非重整化性质，证明它们存在往往很容易。Poincaré 超对称与共形对称不对易，两者合起来生成超共形对称群。超共形代数（SCA）由自旋 $\frac32$ 流的个数 $\mathcal{N}$ 刻画（但不唯一）。已知有 CFT 容许 $\mathcal{N}=1$、$2$、$4$ 的全纯与反全纯 SCA。

$\mathcal{N}=4$ SCA 分两类，性质很不一样：小的 $\mathcal{N}=4$ SCA 只含一个 $SU(2)$“R 对称性”流代数；大的 $\mathcal{N}=4$ SCA 含 $SU(2)\times SU(2)\times U(1)$ 流代数。

SCA 是对玻色 Virasoro 代数的扩充，因此超共形场论（SCFT）可以看作带额外高自旋对称性的普通玻色 CFT，但有一个重要的微妙之处：SCA 含半整数自旋流，这似乎破坏柱面上的动量量子化，也破坏模不变性。要让理论自洽，有两条路：要么稍微放宽模不变性的通常概念，允许谱中出现半整数自旋算符；要么坚持整数自旋条件，把半整数自旋算符从谱中投影掉。后者称为 Gliozzi--Scherk--Olive（GSO）投影 \cite{r26}，做法如下。半整数自旋算符是反对易场，具有奇费米子数 $F$；$F$ 是模 $2$ 意义下的整数。投影到整数自旋算符，等于投影到 $(-)^F$ 不变的态。要与模不变性相容，还必须纳入关于 $(-)^F$ 的扭曲扇区态。在 GSO 投影后的理论中，半整数自旋的超共形流严格说并不在谱内，但仍可用它们按 SCA 的表示来组织谱。

下面讨论 SCA 的几个例子（远非穷尽）以及实现这些对称性的 CFT。先从 $\mathcal{N}=1$ SCA 开始：除中心荷为 $c$ 的能动张量 $T(z)$ 外，它只含一个自旋 $\frac32$ 的反对易场 $G(z)$，且 $G(z)$ 是 Virasoro 原初场。唯一新的独立 OPE 关系是
\begin{equation}
  \label{eq:4-28}
  G(z)G(0)\sim\frac{2c}{3z^3}+\frac{2}{z}T(0).
\end{equation}
作用在原点处的态上时，用 Laurent 级数
\begin{equation}
  \label{eq:4-29}
  G(z)=\sum_{r\in\mathbb{Z}+\nu}\frac{G_r}{z^{r+\frac32}}
\end{equation}
更方便。

% [Tutor Note: $\nu=\frac12$ 时 $r\in\mathbb{Z}+\frac12$，对应 NS 扇区；$\nu=0$ 时 $r\in\mathbb{Z}$，对应 R 扇区。作者用同一个 $\nu$ 统一写出两种展开，这是核对 (4.36)–(4.38) 中 $q$ 的幂次与模指标时最容易出错的地方：NS 扇区态按半整数模指标计数，R 扇区态按整数模指标计数。]

其中在无支割线时 $\nu=\frac12$。此时称原点处的算符处于 Neveu--Schwarz（NS）扇区。特别地，单位算符属于 NS 扇区。GSO 投影要求存在扭曲扇区算符，并在其上附着一条与 $(-)^F$ 相伴的拓扑线算符。这条拓扑线算符造成一条支割线，结果 $G(z)$ 的 Laurent 级数 \eqref{eq:4-29} 中 $\nu=0$。这类扭曲扇区算符称为 Ramond（R）扇区算符。

$\mathcal{N}=(1,1)$ SCFT 的一些基本例子，是中心荷 $c=\frac32-\frac{12}{(k+2)(k+4)}$（$k=1,2,\cdots$）的 $\mathcal{N}=1$ 极小模型。有趣的是，GSO 投影后的 $\mathcal{N}=1$ 极小模型也可实现为玻色陪集模型 $(SU(2)_k\times SU(2)_2)/SU(2)_{k+2}$。另一些例子包括自由 $c=\frac32$ 理论，它由一个紧致玻色子和一对互为手征、反手征的自由费米子 $\psi,\tilde\psi$ 组成；还有 $\mathcal{N}=1$ Liouville CFT：在拉氏量层面，它可以看作线性膨胀子 CFT 加上自由费米子 $\psi,\tilde\psi$，并在领头阶受到形式为 $\psi\tilde\psi e^{\beta\phi}$ 的边缘算符形变。

接着考虑 $\mathcal{N}=2$ SCA。它有一对自旋 $\frac32$ 的反对易场 $G^\pm(z)$ 和一个 $U(1)$ R 流 $J(z)$。可把 $J(z)$ 归一化，使其 OPE 取形式
\begin{equation}
  \label{eq:4-30}
  J(z)J(0)\sim\frac{c}{3z^2},
\end{equation}
此时 $G^\pm$ 相对 $J(z)$ 带电荷 $\pm1$。$G^+G^+$ 与 $G^-G^-$ 的 OPE 非奇异。唯一新的非平凡 OPE 关系是
\begin{equation}
  \label{eq:4-31}
  G^+(z)G^-(0)\sim\frac{2c}{3z^3}+\frac{2}{z^2}J(0)+\frac{2}{z}T(0)+\frac{1}{z}\partial J(0).
\end{equation}

% [Tutor Note: (4.31) 末项系数极易按记忆写成 $2$。作者采用 $\{G^+_r,G^-_s\}=2L_{r+s}+(r-s)J_{r+s}+\frac{c}{3}(r^2-\frac14)\delta_{r+s,0}$ 与 $J(z)J(0)\sim\frac{c}{3z^2}$ 的归一化：用 $[J_2,\cdot]$ 作用于 (4.31) 两边（$[J_2,G^\pm(z)]=\pm z^2G^\pm(z)$，$[J_2,\partial J(0)]=\frac{2c}{3}$），即定出该系数为 $1$。]

用 Laurent 模式 $G^\pm_r,L_n,J_n$ 表示，$\mathcal{N}=2$ SCA 有一个外自同构，称为谱流：
\begin{equation}
  \label{eq:4-32}
  G^\pm_r\to G^\pm_{r\pm\eta},\qquad J_n\to J_n+\eta\frac{c}{3}\delta_{n,0},\qquad L_n\to L_n+\eta J_n+\eta^2\frac{c}{6}\delta_{n,0},
\end{equation}
其中 $\eta\in\mathbb{Z}$。也可以把谱流延拓到半整数 $\eta$，它连接 NS 与 R 扇区。不过，谱流不必是 SCFT 的对称性。

某些“好”的 $\mathcal{N}=2$ SCFT 确实具有谱流对称性。一个例子是 Calabi--Yau 靶空间上的 $(2,2)$ NLSM。此时可按 $\mathcal{N}=2$ SCA 加上谱流对称性（$\eta$ 取整数或半整数）生成的更大对称性来组织算符谱。这称为“扩展 $\mathcal{N}=2$ SCA”\cite{r27,r28}。

% [Tutor Note: 谱流 (4.32) 只是 $\mathcal{N}=2$ SCA 的外自同构，一般不保持谱；在 Calabi--Yau 上的 $(2,2)$ NLSM 里它成为真正的对称性，原因是靶空间的复结构给出额外的整体 $U(1)$，使整数（乃至半整数）$\eta$ 的谱流把谱映回自身。正因如此，下文的 Calabi--Yau 模型才能用扩展 $\mathcal{N}=2$ SCA 的表示论描述 BPS 谱，而不必处理一般 $\mathcal{N}=2$ SCA 的种种缩短表示。]

$\mathcal{N}=2$ SCFT 的其他例子包括 $\mathcal{N}=2$ 极小模型，$c=\frac{3k}{k+2}$，$k=1,2,\cdots$。注意它们的中心荷与 $SU(2)_k$ WZW 模型相同，但并不是同一个 CFT（特别是算符谱不同）。另一个非平凡例子是 $\mathcal{N}=2$ Liouville CFT \cite{r29}：在拉氏量层面，它可以看作紧致玻色子加线性膨胀子再加一对费米子，并在领头阶受到形式为 $\psi_+\tilde\psi_+e^{-\beta\varphi}+\psi_-\tilde\psi_-e^{-\beta\bar\varphi}$ 的边缘算符形变。

最后讨论小的 $\mathcal{N}=4$ SCA。它有四个自旋 $\frac32$ 的流 $G^{\alpha A}(z)$，其中 $\alpha$ 是关于 $SU(2)_k$ R 流 $J_i(z)$ 的二重态指标，$c=6k$；$A$ 是一个额外指标，在外自同构 $SU(2)_{\mathrm{out}}$ 下变换，后者不属于超共形对称性，一般也不是 SCFT 的对称性。注意 $SU(2)$ 能级 $k$ 是整数，故 $c$ 必须是 $6$ 的整数倍。唯一新的非平凡 OPE 关系是
\begin{equation}
  \label{eq:4-33}
  G^{\alpha A}(z)G^{\beta B}(0)\sim\varepsilon^{AB}\Big[\frac{2c}{3z^3}\varepsilon^{\alpha\beta}
  +\sigma_i^{\alpha\beta}\frac{2}{z^2}J_i(0)+\varepsilon^{\alpha\beta}\frac{2}{z}T(0)
  +\sigma_i^{\alpha\beta}\frac{1}{z}\partial J_i(0)\Big].
\end{equation}

% [Tutor Note: (4.33) 中 $\sigma_i^{\alpha\beta}$（$i=1,2,3$）是 $SU(2)$ 的 Pauli 矩阵，$\varepsilon^{\alpha\beta}$、$\varepsilon^{AB}$ 是反对称张量；方括号内四项分别是中心项、R 流、能动张量与 R 流的导数，指标缩并方式决定整体符号。小 $\mathcal{N}=4$ 的 $c=6k$ 与式中的 $2c/3=4k$ 一致，可用它快速检验系数。]

$\mathcal{N}=4$ SCA 也容许谱流映射
\begin{equation}
  \label{eq:4-34}
  G^{\pm A}_r\to G^{\pm A}_{r\pm\eta},\qquad J^3_n\to J^3_n+\eta k\delta_{n,0},\qquad
  J^\pm_n\to J^\pm_{n\pm2\eta},\qquad L_n\to L_n+2\eta J^3_n+\eta^2k\delta_{n,0},
\end{equation}
其中 $\eta$ 为整数（半整数 $\eta$ 会联系 NS 与 R 扇区）。这里相当不平凡的是，谱流本身是 $\mathcal{N}=4$ SCA 的一部分；换句话说，与 $\mathcal{N}=2$ SCA 不同，小 $\mathcal{N}=4$ SCA 自动是“扩展”的。这可通过研究 $\mathcal{N}=4$ SCA 的真空特征标并看出它在谱流下不变来证明。

% [Tutor Note: 对比 (4.32) 与 (4.34)：$\mathcal{N}=2$ 时 $J_n$ 的平移量是 $\eta c/3$，$\mathcal{N}=4$ 时 $J^3_n$ 的平移量是 $\eta k=\eta c/6$，而 $L_n$ 的二次项分别是 $\eta^2c/6$ 与 $\eta^2k=\eta^2c/6$，两处一致。注意 $J^\pm_n$ 的指标平移是 $2\eta$、而 (4.32) 中 $G^\pm_r$ 的平移是 $\pm\eta$：两式对谱流参数 $\eta$ 的归一化约定不同，合读时不要按同一 $\eta$ 数值代换。]

小 $\mathcal{N}=4$ SCA 的表示论相对简单：幺正表示由权 $h$ 与 $SU(2)_R$ 自旋 $\ell$ 标记，$\ell$ 取有限范围 \cite{r30}。在 NS 扇区，非 BPS 原初态满足 $h>\ell$，$\ell=0,\frac12,\cdots,\frac{k-1}{2}$；BPS 原初态满足 $h=\ell$，$\ell=0,\frac12,\cdots,\frac{k}{2}$。

容许小 $\mathcal{N}=4$ SCA 的 SCFT 的一类非平凡例子，是 hyperKähler 流形上的超对称 NLSM。另有一族 $\mathcal{N}=4$ Liouville 理论 \cite{r31}，中心荷 $c=6k$，可由线性膨胀子、四个自由费米子与 $SU(2)_{k-1}$ WZW 模型组成的 CFT 经边缘形变构造。

\subsection{Calabi--Yau 模型}

Calabi--Yau $d$ 重流形上的超对称 NLSM 配上一族合适的度规和可能的 $B$ 场（它们解出经 $\alpha'$ 修正的 Ricci 平坦条件），定义了一族 $c=3d$、具有扩展 $\mathcal{N}=2$ SCA 的 $\mathcal{N}=(2,2)$ SCFT。特别地，全纯 NS 扇区上的 $\eta=1$ 谱流把单位算符（NS 扇区真空态）映成权 $h=\frac d2$、$U(1)_R$ 荷为 $d$ 的超共形原初算符 $\mathcal{O}_\Omega$。并且 $\mathcal{O}_\Omega$ 被 $G^+_{-\frac12}$ 湮灭，是一个手征原初算符。它在 NLSM 描述里有自然的解释：存在全纯 $d$-形式 $\Omega=\frac{1}{d!}\Omega_{i_1\cdots i_d}dz^{i_1}\wedge\cdots\wedge dz^{i_d}$，由此可构造算符
\begin{equation}
  \label{eq:4-35}
  \mathcal{O}_\Omega\sim\Omega_{i_1\cdots i_d}\psi^{i_1}\cdots\psi^{i_d}.
\end{equation}
在相互作用 NLSM 中，右端先验地需要正规化；但实际上由于 $\Omega$ 的全纯性，这样的算符不会被重整化（注意 $\bar z^{\bar i}$ 与 $\bar\psi^{\bar i}$ 都不出现）。

% [Tutor Note: $\mathcal{O}_\Omega$ 是 $(c,c)$ 环中的手征原初算符：它是 $(d,0)$ 形式与手征费米子的缩并，全纯性使它不受 $\bar z$、$\bar\psi$ 方向的短距离奇异性影响，因而不被重整化。这类算符在拓扑 twist 后成为拓扑弦的可观测量，也是下文计数 BPS 态的出发点。]

到目前为止，我们对作为 SCFT 的 Calabi--Yau 模型的认识，大体限于精确边缘形变的模空间、BPS 算符谱，以及由共形自举得到的关于非 BPS 算符谱的一些粗糙限制。例外是模空间中的一些特殊点（称为 Gepner 点），那里的 SCFT 恰好是有理的，并可描述为极小模型乘积的轨形。

下面简述算符谱。一般 $\mathcal{N}=2$ SCA 的表示论相当精细 \cite{r32,r33}，包含多种不同的缩短表示；扩展 $\mathcal{N}=2$ SCA 的表示论则相对简单 \cite{r27,r28}。在全纯与反全纯两侧各自独立地，存在非 BPS（有质量）表示与 BPS（无质量）表示。非 BPS 表示含一个超共形原初态，其 R 荷 $Q$ 取整数值，范围是 $Q=0,\pm1,\pm2,\cdots,\pm(\lfloor d/2\rfloor-1)$，$d$ 为奇数时还有 $\frac{d-1}{2}$；它满足 BPS 界 $h>\frac{|Q|}{2}$。非 BPS 表示的全部权与荷内容由下面的 NS 特征标概括：
\begin{equation}
  \label{eq:4-36}
  \begin{aligned}
    ch^{NS}_{h,Q}(q,y)&\equiv\mathrm{Tr}_{\text{non-BPS}}\,q^{L_0-\frac{c}{24}}y^{J_0}\\
    &=q^{h-\frac{c}{24}}\sum_{m=-\infty}^{\infty}q^{\frac{d-1}{2}m^2+mQ}y^{(d-1)m+Q}
    \prod_{n=1}^{\infty}\frac{(1+yq^{n-\frac12})(1+y^{-1}q^{n-\frac12})}{(1-q^n)^2}
  \end{aligned}
\end{equation}

另一方面，BPS 表示含一个手征或反手征原初态，其 R 荷 $Q$ 取整数值，范围是 $Q=0,\pm1,\pm2,\cdots,\pm\big\lfloor\frac{d-1}{2}\big\rfloor$，$d$ 为偶数时还有 $\frac d2$；它饱和 BPS 界 $h=\frac{|Q|}{2}$。相应的 NS 特征标 $\chi^{NS}_Q(q,y)\equiv\mathrm{Tr}_{\text{BPS}}\,q^{L_0-\frac c{24}}y^{J_0}$ 由下式给出：
\begin{equation}
  \label{eq:4-37}
  \begin{aligned}
    \chi^{NS}_Q(q,y)&=q^{\frac Q2-\frac c{24}}\sum_{m\in\mathbb{Z}}\frac{q^{\frac{d-1}{2}m^2+mQ}y^{(d-1)m+Q}}{1+yq^{m+\frac12}}
      \prod_{n=1}^{\infty}\frac{(1+yq^{n-\frac12})(1+y^{-1}q^{n-\frac12})}{(1-q^n)^2},\quad Q>0,\\
    \chi^{NS}_Q(q,y)&=\chi^{NS}_{-Q}(q,y^{-1}),\quad Q<0,\\
    \chi^{NS}_0(q,y)&=q^{-\frac c{24}}\sum_{m\in\mathbb{Z}}\frac{(q^{-\frac12}-q^{\frac12})q^{\frac{d-1}{2}m^2+m}y^{(d-1)m+1}}{(1+yq^{m-\frac12})(1+yq^{m+\frac12})}
      \prod_{n=1}^{\infty}\frac{(1+yq^{n-\frac12})(1+y^{-1}q^{n-\frac12})}{(1-q^n)^2}.
  \end{aligned}
\end{equation}

% [Tutor Note: (4.36) 与 (4.37) 的差别只在分母上的 $(1+yq^{m+\frac12})$ 型因子：非 BPS 表示不缩短，直接对谱流轨道求和；BPS 表示被 $h=\frac{|Q|}{2}$ 缩短，$Q>0$ 时留下 $1+yq^{m+\frac12}$ 的零点，$Q<0$ 由 $\chi^{NS}_Q(q,y)=\chi^{NS}_{-Q}(q,y^{-1})$ 镜像给出，$Q=0$ 则靠 $(q^{-\frac12}-q^{\frac12})$ 抵消分母的零点。核对 (4.36)–(4.37) 时最容易丢掉这些因子。]

同时在两侧都属于 BPS 表示的算符，笼统称为 $\frac12$-BPS；一侧 BPS、另一侧非 BPS 的算符，则称为 $\frac14$-BPS。

有一个超对称指数对约束 $\frac12$-BPS 与 $\frac14$-BPS 算符的谱非常有用，即椭圆亏格 \cite{r34}。定义这个指数必须在基态保持超对称的 Ramond 扇区里做。R 扇区特征标与 NS 扇区特征标由 $\eta=\frac12$ 的谱流相联系，代价是原初权与 R 荷的平移：
\begin{equation}
  \label{eq:4-38}
  ch^{R}_{h,Q}(q,y)=ch^{NS}_{h+\frac c{24}-\frac Q2,\frac c6-Q}\big(q,q^{\frac12}y\big),\qquad
  \chi^{R}_Q(q,y)=\chi^{NS}_{\frac c6-Q}\big(q,q^{\frac12}y\big).
\end{equation}

椭圆亏格定义为 (R,R) 扇区中的特征标，形如
\begin{equation}
  \label{eq:4-39}
  \chi(\tau,z)=\mathrm{Tr}_{RR}\,q^{L_0-\frac c{24}}y^{J_0}\bar q^{\bar L_0-\frac c{24}}e^{\pi i\bar J_0},
\end{equation}
其中 $q=e^{2\pi i\tau}$，$y=e^{2\pi iz}$。注意 $\bar\tau$ 依赖会消失，原因是反全纯 R 荷 $\bar J_0$ 为偶与为奇的激发态之间发生相消。由于（整数 $\eta$ 的）谱流对称性，$\chi$ 在 $z\to z+1$ 与 $z\to z+\tau$ 下不变。此外，$\chi$ 具有 $PSL(2,\mathbb{Z})$ 变换性质
\begin{equation}
  \label{eq:4-40}
  \chi(\tau+1,z)=\chi(\tau,z),\qquad
  \chi(-1/\tau,z/\tau)=e^{2\pi i\frac d2\frac{z^2}{\tau}}\chi(\tau,z).
\end{equation}

% [Tutor Note: (4.39) 中 $e^{\pi i\bar J_0}$ 与 $\bar q^{\bar L_0-\frac c{24}}$ 合起来实现 $(-)^{\bar J_0}$ 加权，反全纯 R 荷奇偶的态两两相消，这就是 $\bar\tau$ 依赖消失的原因；能存活下来的正是被 $\bar G^-_{-\frac12}$、$\bar G^+_{-\frac12}$ 湮灭的 BPS 态。指标 $d/2$ 等于 $\eta=\frac12$ 谱流对 $J_0$ 的平移量 $c/6$，这也是 (4.40) 中 $e^{2\pi i\frac d2\frac{z^2}{\tau}}$ 因子的来源。]

也就是说，$\chi(\tau,z)$ 是权 $0$、指标 $\frac d2$ 的弱 Jacobi 形式。整数指标的弱 Jacobi 形式空间由四个基本 Jacobi 形式 $E_4,E_6,\phi_{0,1},\phi_{-2,1}$ 的多项式张成，其中 $E_4,E_6$ 是 Eisenstein 级数，$\phi_{0,1},\phi_{-2,1}$ 是两个指标 $1$ 的弱 Jacobi 形式。对半整数指标，还可额外使用另外两个基本弱 Jacobi 形式 $\phi_{0,\frac32}$ 与 $\phi_{-1,\frac12}$。

$d=2$ 时（例子是靶空间为 K3 的 SCFT），椭圆亏格必须正比于 $\phi_{0,1}$。把 Witten 指数与 K3 的 Euler 示性数匹配，可得 \cite{r35}
\begin{equation}
  \label{eq:4-41}
  \chi_{\mathrm{K3}}(\tau,z)=2\phi_{0,1}(\tau,z)=2y+20+2y^{-1}+(20y^2-128y+216-128y^{-1}+20y^{-2})q+\cdots.
\end{equation}
在假定小 $\mathcal{N}=4$ SCA 之外没有额外守恒高自旋流的前提下，这可以用来确定 K3 CFT 中所有 $\frac12$-BPS 与 $\frac14$-BPS 态的简并度。

% [Tutor Note: 系数 $2$ 由 Witten 指数定出：$\phi_{0,1}$ 在 $y=1$ 处给 $12$，故 $2\phi_{0,1}$ 给 $24=\chi(\mathrm{K3})$。$q^0$ 层的 $2y+20+2y^{-1}$ 就是 K3 上 $\frac12$-BPS 态的计数；能如此干净地逐层读出简并度，前提正是小 $\mathcal{N}=4$ 之外没有额外高自旋流。]

$d=3$ 时（含 Calabi--Yau 三重流形模型），椭圆亏格必须正比于 $\phi_{0,\frac32}$。把 Witten 指数与 Euler 示性数匹配，得到
\begin{equation}
  \label{eq:4-42}
  \chi_{\mathrm{CY}_3}(\tau,z)=(h^{1,1}-h^{2,1})\phi_{0,\frac32}(\tau,z).
\end{equation}
于是可以把 $\phi_{0,\frac32}(\tau,z)$ 按 RR 扇区的扩展 $\mathcal{N}=2$ 特征标分解。结果很简单：$\phi_{0,\frac32}=\chi^R_1-\chi^R_{-1}$。这个结果对 BPS 谱有有趣的推论：虽然有（扩展 $\mathcal{N}=2$）$\frac12$-BPS 多重态对应 SCFT 的形变模，但 $\frac14$-BPS 多重态要么含有额外守恒的高自旋流（预期在共形流形的一般点上不存在），要么必须两两配对，从而可能在边缘形变下被提升为非 BPS 多重态。

% [Tutor Note: 分解式 $\phi_{0,\frac32}=\chi^R_1-\chi^R_{-1}$ 中，$\chi^R_1$ 与 $\chi^R_{-1}$ 分别是 $Q=\pm1$ 的 BPS 表示，正负号来自 R 扇区中手征与反手征表示的符号。系数只有 $\pm1$，故 $\frac14$-BPS 态无法单独出现，只能配对或伴生高自旋流，这正是这段推理的关键。]

\subsection{精确边缘形变}

当 CFT 容许一个权 $h=\tilde h=1$ 的边缘原初算符 $\mathcal{O}$ 时，可以在关联函数中插入 $\exp[-\lambda\int d^2z\,\mathcal{O}(z,\bar z)]$ 来形变理论。对有拉氏量描述的理论，这相当于把拉氏量密度用算符 $\mathcal{O}(z,\bar z)$ 形变。共形不变性在形变参数 $\lambda$ 的一阶保持；更高阶一般不保持，因为重合算符的正规化会破坏共形对称性。例如，若 $\mathcal{O}(z,\bar z)\mathcal{O}(0)$ 的 OPE 含另一个边缘原初算符，就会在 $\lambda^2$ 阶生成非平凡的 beta 函数。

% [Tutor Note: “正规化”指重合点 $z\to0$ 处 $\mathcal{O}(z,\bar z)\mathcal{O}(0)$ 的短距离奇异性必须截断；若 OPE 中出现另一边缘算符 $\mathcal{O}'$，重整化会生成正比于 $\lambda^2\int\mathcal{O}'$ 的反项，beta 函数因此非零。注意作者把形变写成 $\exp[-\lambda\int]$，与部分文献用 $g\int$ 的符号约定相反，比较公式时留意。]

理解更高阶共形不变性障碍的一个有用办法，是把该 CFT 看作紧致化玻色弦世界面理论的一部分，连同自由玻色子 $X^\mu$ 与 $b,c$ 鬼场。一个边缘原初算符给出 Minkowski 时空（由 $X^\mu$ 参数化）中无质量标量场的顶点算符 $\mathcal{V}_i=\mathcal{O}e^{ik_i\cdot X}$。共形不变性的第 $n$ 阶障碍，由 $n$ 点弦树图散射振幅在 $k_i\to0$ 同时极限下的值刻画：
\begin{equation}
  \label{eq:4-43}
  \Big\langle\Big[\prod_{i=1}^{n-3}\int d^2z_i\,\mathcal{V}_i(z_i,\bar z_i)\Big]
  \mathcal{V}_{n-2}(0)\mathcal{V}_{n-1}(1)\mathcal{V}'_n(\infty)\Big\rangle,
\end{equation}
其中 $z_i$ 积分可通过对 $k_i$ 的解析延拓正规化。

% [Tutor Note: (4.43) 是固定 $SL(2,\mathbb{C})$ 后剩下的 $n-3$ 个重积分；$\mathcal{V}'_n(\infty)$ 表示在无穷远处按共形权插入顶点算符。所谓“零动量极限的障碍”是指：若该振幅在 $k_i\to0$ 时不消失，有效作用量就会生成相应局域反项，破坏共形不变性。$d^2z_i$ 积分在 $k_i$ 非零时发散，作者用 $k_i$ 的解析延拓来定义它，这是 Polchinski 共形微扰论的标准做法。]

精确边缘形变是导致一族连续 CFT 的边缘形变，由形变参数 $\lambda$ 参数化。这要求 \eqref{eq:4-43} 的零动量极限对所有 $n$ 都消失。这个形变空间上有自然的度规概念，由边缘算符的两点函数给出，称为 Zamolodchikov 度规 \cite{r36}。因此，精确边缘形变的这个模空间是 Riemann 流形，有时称为共形流形。基本例子是基于 Narain 格 $\Gamma_{n_L,n_R}$ 的自由玻色子 CFT，它有 $n_Ln_R$ 维共形流形，同构于 $\mathrm{Aut}(\Gamma_{n_L,n_R})\backslash O(n_L,n_R)/(O(n_L)\times O(n_R))$。更一般地，只要 CFT 有一组 $U(1)^{n_L}$ 全纯流 $j^a_z$ 与 $U(1)^{n_R}$ 反全纯流 $\bar j^b_{\bar z}$，则形变 $j^a_z\bar j^b_{\bar z}$ 都精确边缘。

% [Tutor Note: Narain 情形的模空间维数是 $n_Ln_R$，不是 $n_L^2+n_R^2$：商空间里左乘 $\mathrm{Aut}(\Gamma_{n_L,n_R})$ 是格自同构，右除的 $O(n_L)\times O(n_R)$ 分别是左右手征基的转动。$j^a_z\bar j^b_{\bar z}$ 之所以精确边缘，是因为流守恒 $\bar\partial j=0$ 使形变的 OPE 里不会出现会生成 beta 函数的边缘算符。]

现在讨论 $(2,2)$ SCFT 的边缘形变，它们一阶保持全部超共形对称性。这类形变基于权 $(\frac12,\frac12)$ 的（手征，手征）或（手征，反手征）算符，通过形如
\begin{equation}
  \label{eq:4-44}
  \mathcal{O}=G^-_{-\frac12}\bar G^-_{-\frac12}\phi^{+,+}+G^+_{-\frac12}\bar G^+_{-\frac12}\phi^{-,-}
  \quad\text{or}\quad
  G^-_{-\frac12}\bar G^+_{-\frac12}\phi^{+,-}+G^+_{-\frac12}\bar G^-_{-\frac12}\phi^{-,+}
\end{equation}
的后代实现，其中 $\phi^{+,+}$ 表示全纯与反全纯 $U(1)_R$ 荷为 $(1,1)$ 的（手征，手征）原初场，$\phi^{+-}$ 表示 R 荷为 $(1,-1)$ 的（手征，反手征）原初场。要看出这保持超共形对称性，需证明对任意全纯函数 $v(z)$，超共形流 $v(z)G^+(z)$ 可移过 $\mathcal{O}(w,\bar w)$ 而只差一个全导数。在（手征，手征）情形验证如下：
\begin{equation}
  \label{eq:4-45}
  \begin{aligned}
  &\oint\frac{dz}{2\pi i}v(z)G^+(z)\mathcal{O}(w,\bar w)\\
  &\quad=\oint\frac{dz}{2\pi i}v(z)\Big((z-w)^{-2}G^+_{\frac12}+(z-w)^{-1}G^+_{-\frac12}\Big)
    G^-_{-\frac12}\bar G^-_{-\frac12}\phi^{+,+}(w,\bar w)\\
  &\quad=\big[\partial v(w)(2L_0+J_0)+v(w)2L_{-1}\big]\bar G^-_{-\frac12}\phi^{+,+}(w,\bar w)\\
  &\quad=2\partial_w\Big[v(w)\bar G^-_{-\frac12}\phi^{+,+}(w,\bar w)\Big],
  \end{aligned}
\end{equation}
其中用到假定：$G^+_r\phi^{+,+}=0$ 对所有 $r\ge-\frac12$ 成立。

% [Tutor Note: (4.45) 第一行把 $G^+$ 的围道积分写成 $\{G^+_{\frac12},G^-_{-\frac12}\}(z-w)^{-2}+\{G^+_{-\frac12},G^-_{-\frac12}\}(z-w)^{-1}$ 两项，第二步用到 $\{G^+_{\frac12},G^-_{-\frac12}\}=2L_0+J_0$ 与 $\{G^+_{-\frac12},G^-_{-\frac12}\}=2L_{-1}$；后者没有 $J_{-1}$ 项，因为 $r-s=0$。$G^+_r\phi^{+,+}=0$（$r\ge-\frac12$）消掉其余项。核对符号时注意 $G^+$ 与 $G^-$ 反对易。]

$(2,2)$ SCFT 的一个重要性质是：边缘形变必定精确边缘 \cite{r37,r38}，由此得到大量（超）共形流形的例子。基本理由是：当我们用例如 $\lambda\int d^2z\,G^-_{-\frac12}\bar G^-_{-\frac12}\phi^{+,+}+c.c.$ 形变理论，并以超对称方式正规化时，$\lambda$ 只能全纯地重整化。这来自把形变写成超空间积分 $\int d^2z\,d^2\theta^+\,\lambda\Phi^{+,+}(z,\bar z,\theta,\bar\theta)+c.c.$，其中 $\Phi^{+,+}$ 是手征超场，最低分量为 $\phi^{+,+}$；接着把 $\lambda$ 提升为手征超参数（也称“spurion”），并考虑 RG 下可能生成的所有局域超空间积分。于是，由 $\phi^{+,+}$ 之间 OPE 非奇异可知，$\lambda$ 根本不能被重整化。

% [Tutor Note: “spurion”技巧是把耦合常数提升为带指定超对称变换性质的外部超场：这里 $\lambda$ 成为手征超参数，只允许依赖 $\theta$ 而不依赖 $\bar\theta$。因此任何重整化只能给出全纯的 $\lambda$ 修正；而 $\phi^{+,+}$ 之间的 OPE 不含其他手征原初场，故不存在可生成的反项，$\lambda$ 保持不动。]

一个有趣例子是 D1--D5 CFT \cite{r39,r40,r41}，它展示了共形流形上可能出现的某些丰富现象。它是具有小 $\mathcal{N}=(4,4)$ SCA、中心荷 $c=6k$（$k$ 为正整数）的 SCFT。在共形流形的某些极限下，D1--D4 CFT 可由超对称 NLSM 描述，其靶空间是 K3 上 $U(Q_5)$ 自对偶 Yang--Mills 方程解（瞬子数为 $Q_1+Q_5$，对应 $k=Q_1Q_5+1$）的模空间，或 $T^4$ 上瞬子数为 $Q_1$（对应 $k=Q_1Q_5$）的解的模空间。特别地，$Q_5=1$ 的情形对应轨形 SCFT $\mathrm{Sym}^{Q_1Q_5+1}(\mathrm{K3})$ 或 $\mathrm{Sym}^{Q_1Q_5}(T^4)$。当 $Q_1$ 与 $Q_5$ 互素时，预期 D1--D5 CFT 是紧致的。基于弦对偶的论证表明：若 $(Q_1,Q_5)$ 与 $(Q_1',Q_5')$ 都与 $Q_1Q_5=Q_1'Q_5'$ 互素，则由 $(Q_1,Q_5)$ 与 $(Q_1',Q_5')$ 标记的 D1--D5 CFT 位于同一共形流形上。这意味着，看似非常不同的靶空间几何上的共形 NLSM，通过强耦合区域（靶空间变得高度弯曲）中的精确边缘形变以某种方式相互联系。

% [Tutor Note: 原文此处的“D1--D4 CFT”按上下文应为“D1--D5 CFT”之笔误：本段其余部分、对 $Q_5$ 的依赖以及 $Q_5=1$ 时给出的 $\mathrm{Sym}^{Q_1Q_5+1}(\mathrm{K3})$ 都指 D1--D5 系统。译文照原文排作 D1--D4，核对时按 D1--D5 理解。]

此外，模空间中还有些点，使 D1--D5 CFT 在某个维数间隙之上出现连续谱（这些连续态是在整体 $AdS_3$ 中以有限能量代价跑到无穷远的长闭弦的全息对偶）。

\subsection{重整化群流}

受相关算符微扰的 CFT 会在重整化群下流向中心荷 $c$ 更小的另一个 CFT。这是探索 CFT 之间关系、构造潜在新 CFT 的重要工具。从 UV CFT 出发，红外 RG 不动点一般超出共形微扰论的能力；因此，除非已知 RG 流可积或被对称性（通常是超对称）保护，否则精确确定红外 CFT 相当困难。这里简短提及二维中几类可控的 RG 流。

第一类是可积的 RG 流 \cite{r42,r43,r44,r45,r46,r47}。最简单的例子是 IR 流向有质量理论，其中含有限种基本粒子和一个因子化的 $S$ 矩阵。此时，UV CFT 经可积相关微扰得到的 QFT 完全由 IR 粒子谱与两体 $S$ 矩阵元决定。由这些数据，原则上可用热力学 Bethe 拟设（TBA）确定圆上 QFT 的谱，并从半径趋于零的极限恢复 UV CFT。此外，局域算符的关联函数可由形状因子构造，而形状因子受一组解析性约束以及由因子化 $S$ 矩阵决定的交叉对称性支配 \cite{r48,r49}。

% [Tutor Note: 可积流的关键是无穷多守恒流把 $S$ 矩阵完全固定为因子化形式（Yang--Baxter 方程的解），故两体振幅足以确定整个理论；TBA 给出有限温度或有限体积配分函数的积分方程，$R\to0$ 时其渐近行为编码 UV 中心荷，这是用 TBA 反推 UV CFT 的常用手段。]

第二类是 Landau--Ginzburg（LG）模型的 RG 流。在玻色情形，已知非紧致自由玻色子 $\phi$ 的 CFT 受形式为 $\phi^{2n}$ 的算符形变后，会生成一条精细调节的 RG 流，流向中心荷为 $c=1-\frac{6}{(n+1)(n+2)}$ 的极小模型。$n$ 大时可用共形微扰论检验；$n$ 小时可用截断共形空间方法，直接对角化形变后的 Hamiltonian，并用 Rayleigh--Ritz 方法 \cite{r50,r51,r52}。

% [Tutor Note: $\phi^{2n}$ 形变的“精细调节”指耦合常数必须取特定值才落到极小模型；一般耦合下 RG 流奔向有质量或平庸的红外。$c=1-\frac{6}{(n+1)(n+2)}$ 正是极小模型 $M(n+1,n+2)$ 的中心荷，$n=2$ 对应 Ising（$c=\frac12$）。截断共形空间方法把希尔伯特空间截断到有限个 Virasoro 表示再数值对角化，误差随截断能标指数减小。]

控制得更好的是 LG 模型的 $(2,2)$ 超对称版本 \cite{r53}：它由一组手征超场 $X_i$ 的拉氏量描述，带超势 $W(X_i)$。只要一开始没有质量项（即场中的二次项），超势的非重整化性质就保证红外不动点非平凡。LG 模型对确定红外 CFT 的共形流形以及 BPS 算符的结构常数特别有用。

第三个例子是一类“短”RG 流，类似三维临界 $O(N)$ 矢量模型。先从 $N$ 份 $c=\frac45$ 临界三态 Potts 模型（也称 $(A_4,D_4)$ 极小模型）的张量积出发。三态 Potts 模型含非零自旋的 Virasoro 原初场，它们由 $c=\frac45$ Virasoro 代数的不同全纯与反全纯退化原初场构造。它有 $\mathbb{Z}_3$ 对称性，在此对称性下，权 $(\frac1{15},\frac1{15})$ 的原初场 $\sigma$ 及其共轭 $\sigma^*$ 分别带相位 $e^{2\pi i/3}$ 与 $e^{-2\pi i/3}$，而权 $(\frac25,\frac25)$ 的“能量算符” $\varepsilon$ 不带荷。$N$ 重张量积 CFT 只有两个在 $S_N\ltimes\mathbb{Z}_3^N$ 下不变的相关算符：
\begin{equation}
  \label{eq:4-46}
  \sum_{i=1}^{N}\varepsilon_i,\qquad \sum_{i<j}^{N}\varepsilon_i\varepsilon_j.
\end{equation}
前一个会生成有质量的 RG 流。后一个生成流向新不动点的 RG 流，至少在 $N$ 足够大时如此；此时可考虑 $1/N$ 展开，并可视算符 $\sum_{i<j}\varepsilon_i\varepsilon_j$ 为双迹形变。这导致一条短 RG 流，中心荷的改变预期为 $1/N$ 阶。确实，看起来连 $N=3$ 情形给出的 RG 流也非常短 \cite{r54}。

% [Tutor Note: $S_N\ltimes\mathbb{Z}_3^N$ 是 Potts 自旋置换与逐点 $\mathbb{Z}_3$ 旋转的半直积；$\sum_i\varepsilon_i$ 是单迹、$\sum_{i<j}\varepsilon_i\varepsilon_j$ 是双迹，前者对应质量形变，后者只在 $1/N$ 展开里显得小，故流动很“短”。$\sigma$ 带 $\mathbb{Z}_3$ 荷而 $\varepsilon$ 不带荷，这一点决定了不变算符里只能出现这两个组合。]

最后提一下，二维中另一类重要的 RG 流是规范场论，包括规范线性 sigma 模型（GLSM），它们往往在 IR 流向非平凡 CFT（如 Calabi--Yau 模型）\cite{r55}。这些并不属于“CFT 受相关算符微扰”的标准范式，因为规范场论在 UV 不是共形的。\footnote{注意 $c$ 定理完全适用，UV 的 $c$ 函数等于自由物质场的 $c$ 函数。然而 UV 极限并不等于自由物质场理论，原因是对局域算符的规范单态约束，以及理论放在空间圆上时（模不变性所需）的和乐自由度。}最简单的例子是二维纯（自由）$U(1)$ 规范场论，它在圆上的能谱同构于半径为 $R=1/g$（$g$ 为规范耦合）的圆上自由非相对论粒子的能谱。显然，在 $R\to\infty$ 极限下，这个谱不会组织成任何 Virasoro 代数的表示，谱密度也不是标度不变的（与非紧致 CFT 不同）。
